\ifdefined\OTCOMBINED%
\else
\documentclass{otscience}
% ============================================================
% Shared setup for OT Science modular workbook parts
% Place these files in the same folder as your fixed otscience.cls
% ============================================================

\makeatletter
\makeatother
\usepackage{multicol}
\usepackage{longtable}
\usepackage{makecell}
\usepackage{pdflscape}
\usepackage{bm}
\providecommand{\blank}[1][3cm]{\rule{#1}{0.4pt}}
\providecommand{\bigblank}{\rule{7cm}{0.4pt}}
\providecommand{\componentblank}{\blank[2.5cm]}
\providecommand{\R}{\mathbb{R}}
\providecommand{\C}{\mathbb{C}}
\providecommand{\ii}{\mathrm{i}}
\providecommand{\ee}{\mathrm{e}}
\providecommand{\dd}{\mathrm{d}}
\providecommand{\grad}{\nabla}
\providecommand{\divergence}{\nabla\!\cdot}
\providecommand{\curl}{\nabla\!\times}
\providecommand{\lap}{\nabla^2}
\providecommand{\ihat}{\hat{\mathbf{i}}}
\providecommand{\jhat}{\hat{\mathbf{j}}}
\providecommand{\khat}{\hat{\mathbf{k}}}
\providecommand{\er}{\hat{\mathbf{e}}_r}
\providecommand{\etheta}{\hat{\mathbf{e}}_\theta}
\providecommand{\ephi}{\hat{\mathbf{e}}_\phi}
\providecommand{\erho}{\hat{\boldsymbol{\rho}}}
\providecommand{\vA}{\mathbf{A}}
\providecommand{\vB}{\mathbf{B}}
\providecommand{\va}{\mathbf{a}}
\providecommand{\vb}{\mathbf{b}}
\providecommand{\vc}{\mathbf{c}}
\providecommand{\vx}{\mathbf{x}}
\providecommand{\vr}{\mathbf{r}}
\providecommand{\matA}{\mathbf{A}}
\providecommand{\matB}{\mathbf{B}}
\providecommand{\tr}{\operatorname{tr}}
\providecommand{\cof}{\operatorname{cof}}
\providecommand{\dev}{\operatorname{dev}}
\providecommand{\diag}{\operatorname{diag}}
\providecommand{\questionline}[1][5cm]{\rule{#1}{0.4pt}}
\setlength{\parindent}{0pt}
\setlength{\parskip}{4pt}
\renewcommand{\arraystretch}{1.25}

% Fallbacks in case the current otscience.cls has not defined nosplit boxes.
\makeatletter
\@ifundefined{formulaboxnosplit}{\newenvironment{formulaboxnosplit}[1][Formula]{\begin{formulabox}[#1]}{\end{formulabox}}}{}
\@ifundefined{definitionboxnosplit}{\newenvironment{definitionboxnosplit}[1][Definition]{\begin{definitionbox}[#1]}{\end{definitionbox}}}{}
\@ifundefined{summaryboxnosplit}{\newenvironment{summaryboxnosplit}[1][Summary]{\begin{summarybox}[#1]}{\end{summarybox}}}{}
\@ifundefined{warningboxnosplit}{\newenvironment{warningboxnosplit}[1][Warning]{\begin{warningbox}[#1]}{\end{warningbox}}}{}
\@ifundefined{exampleboxnosplit}{\newenvironment{exampleboxnosplit}[1][Example]{\begin{examplebox}[#1]}{\end{examplebox}}}{}
\makeatother

\newenvironment{practicebox}[1][Quick Drills and Practice]{\begin{examplebox}[#1]}{\end{examplebox}}
\newenvironment{practiceboxnosplit}[1][Quick Drills and Practice]{\begin{exampleboxnosplit}[#1]}{\end{exampleboxnosplit}}

\newcommand{\standalonetitle}[3]{%
    \title{\textbf{#1}\\\large #2}
    \author{OT Science Notes}
    \date{#3}
}

\standalonetitle{Covariant and Contravariant Bases, Metrics and General Curvilinear Calculus}{Part 4 of the OT Science Formula Completion Workbook}{\DTMtoday}
\begin{document}
\maketitle
\tableofcontents
\newpage
\fi

% 04_covariant_contravariant_metrics.tex
\section[
    Covariant and Contravariant Bases, Metrics and General Curvilinear Calculus
]{%
    Covariant and Contravariant Bases, Metrics\\
    and General Curvilinear Calculus%
}
\sectionmark{Metrics and Bases}

This part connects coordinate geometry with tensor notation. The key idea is that a vector itself is geometric, but its components depend on the basis used to describe it. Covariant and contravariant components are two related ways of recording the same geometric object.

\subsection{Coordinate Basis and Reciprocal Basis}

Let
\[
    \mathbf{r}=\mathbf{r}(q^1,q^2,q^3).
\]
The covariant basis vectors are tangent to coordinate curves:
\[
    \mathbf{e}_i=\frac{\partial\mathbf{r}}{\partial q^i}.
\]
The contravariant basis vectors form the reciprocal basis:
\[
    \mathbf{e}^i\cdot\mathbf{e}_j=\delta^i_j.
\]

\begin{formulaboxnosplit}[Covariant and Contravariant Bases]
    \[
        \mathbf{e}_i=\frac{\partial\mathbf{r}}{\partial q^i},
        \qquad
        \mathbf{e}^i=\nabla q^i.
    \]

    \[
        \mathbf{e}^i\cdot\mathbf{e}_j=\delta^i_j.
    \]

    \[
        \mathbf{A}=A^i\mathbf{e}_i=A_i\mathbf{e}^i.
    \]

    \[
        A^i=\mathbf{A}\cdot\mathbf{e}^i,
        \qquad
        A_i=\mathbf{A}\cdot\mathbf{e}_i.
    \]
\end{formulaboxnosplit}

\subsection{Metric Tensor}

The metric tensor records dot products between basis vectors. It is the object that turns coordinate changes into physical lengths and angles.

\begin{formulabox}[Metric Tensor and Raising/Lowering]
    \[
        g_{ij}=\mathbf{e}_i\cdot\mathbf{e}_j,
        \qquad
        g^{ij}=\mathbf{e}^i\cdot\mathbf{e}^j.
    \]

    \[
        g^{ik}g_{kj}=\delta^i_j.
    \]

    \[
        A_i=g_{ij}A^j,
        \qquad
        A^i=g^{ij}A_j.
    \]

    \[
        ds^2=g_{ij}\,dq^{i}dq^j.
    \]

    \[
        \mathbf{A}\cdot\mathbf{B}=g_{ij}A^{i}B^j=g^{ij}A_{i}B_j=A_{i}B^i=A^{i}B_i.
    \]
\end{formulabox}

\subsection{Orthogonal Coordinates as a Special Case}

If the coordinate system is orthogonal, the metric tensor is diagonal. This is why scale factors are enough to describe many engineering coordinate systems.

\begin{formulaboxnosplit}[Orthogonal Metric and Scale Factors]
    \[
        g_{ij}=h_i^2\delta_{ij}\quad\text{no sum on }i.
    \]

    \[
        g^{ij}=\frac{1}{h_i^2}\delta^{ij}\quad\text{no sum on }i.
    \]

    \[
        \mathbf{e}_i=h_i\hat{\mathbf{e}}_i,
        \qquad
        \mathbf{e}^i=\frac{1}{h_i}\hat{\mathbf{e}}_i.
    \]

    \[
        A_i=h_i A_{\text{phys},i},
        \qquad
        A^i=\frac{A_{\text{phys},i}}{h_i}.
    \]
\end{formulaboxnosplit}

\subsection{Jacobian, Metric and Volume}

For a coordinate transformation \(x^a=x^a(q^1,q^2,q^3)\), the Jacobian matrix connects small coordinate changes to small physical displacements.

\begin{formulabox}[Jacobian and Volume Element]
    \[
        J^a{}_i=\frac{\partial x^a}{\partial q^i}.
    \]

    \[
        g_{ij}=\frac{\partial x^a}{\partial q^i}\frac{\partial x^a}{\partial q^j}=J^a{}_{i}J^a{}_j.
    \]

    \[
        g=\det(g_{ij}),
        \qquad
        \sqrt{g}=\left|\det(J)\right|.
    \]

    \[
        dV=\sqrt{g}\,dq^1dq^2dq^3.
    \]

    For orthogonal coordinates,
    \[
        \sqrt{g}=h_1h_2h_3.
    \]
\end{formulabox}

\subsection{Gradient, Divergence and Laplacian in General Curvilinear Coordinates}

The following formulas use the metric determinant \(g=\det(g_{ij})\). They are especially useful when coordinates are not necessarily orthogonal.

\begin{formulabox}[General Curvilinear Formulae]
    \[
        \nabla f=(\partial_i f)\mathbf{e}^i.
    \]

    \[
        \nabla\cdot\mathbf{A}=\frac{1}{\sqrt{g}}\frac{\partial}{\partial q^i}\left(\sqrt{g}\,A^i\right).
    \]

    \[
        \nabla^2f=\frac{1}{\sqrt{g}}\frac{\partial}{\partial q^i}\left(\sqrt{g}\,g^{ij}\frac{\partial f}{\partial q^j}\right).
    \]

    For orthogonal coordinates,
    \[
        \nabla^2 f=\frac{1}{h_1h_2h_3}\sum_{i=1}^{3}\frac{\partial}{\partial q^i}\left(\frac{h_1h_2h_3}{h_i^2}\frac{\partial f}{\partial q^i}\right).
    \]
\end{formulabox}

\subsection{Cylindrical and Spherical Laplacians}

\begin{formulabox}[Common Laplacians]
    Cartesian:
    \[
        \nabla^2f=f_{xx}+f_{yy}+f_{zz}.
    \]

    Cylindrical:
    \[
        \nabla^2f=\frac{1}{\rho}\frac{\partial}{\partial\rho}\left(\rho\frac{\partial f}{\partial\rho}\right)+\frac{1}{\rho^2}\frac{\partial^2f}{\partial\phi^2}+\frac{\partial^2f}{\partial z^2}.
    \]

    Spherical:
    \[
        \nabla^2f=\frac{1}{r^2}\frac{\partial}{\partial r}\left(r^2\frac{\partial f}{\partial r}\right)
        +\frac{1}{r^2\sin\theta}\frac{\partial}{\partial\theta}\left(\sin\theta\frac{\partial f}{\partial\theta}\right)
        +\frac{1}{r^2\sin^2\theta}\frac{\partial^2f}{\partial\phi^2}.
    \]
\end{formulabox}

\begin{practicebox}[Quick Drills and Practice: Covariant/Contravariant Components and Metrics]
    \textbf{Level A:\,Conceptual recall}

        \begin{enumerate}
            \item Complete: \(\mathbf{e}_i=\blank[4cm]\).
            \item Complete: \(\mathbf{e}^i=\blank[3cm]\).
            \item Complete: \(\mathbf{e}^i\cdot\mathbf{e}_j=\blank[2cm]\).
            \item Complete: \(g_{ij}=\blank[4cm]\).
            \item Complete: \(g^{ik}g_{kj}=\blank[2cm]\).
            \item Complete: \(A_i=\blank[3cm]\).
            \item Complete: \(A^i=\blank[3cm]\).
            \item Complete: \(ds^2=\blank[4cm]\).
        \end{enumerate}


    \textbf{Level B:\,Metric computations}
    \begin{enumerate}
        \item Compute \(g_{ij}\) for cylindrical coordinates directly from \(\mathbf{r}(\rho,\phi,z)\).
        \item Compute \(g^{ij}\) for cylindrical coordinates.
        \item Compute \(g_{ij}\) for spherical coordinates directly from \(\mathbf{r}(r,\theta,\phi)\).
        \item Compute \(g^{ij}\) for spherical coordinates.
        \item Compute \(\sqrt{g}\) for Cartesian, cylindrical and spherical coordinates.
        \item Show that \(dV=\sqrt{g}\,dq^1dq^2dq^3\) gives the correct cylindrical and spherical volume elements.
    \end{enumerate}

    \textbf{Level C:\,Components}
    \begin{enumerate}
        \item In cylindrical coordinates, if the physical components are \((A_\rho,A_\phi,A_z)=(2,3,4)\), write the contravariant components \(A^i\).
        \item In spherical coordinates, if the physical components are \((A_r,A_\theta,A_\phi)=(1,2,3)\), write \(A^i\) and \(A_i\).
        \item Given \(g_{ij}=\diag(1,r^2,r^2\sin^2\theta)\), lower the index of \(A^i=(r,\sin\theta,1)\).
        \item Given \(g^{ij}=\diag(1,1/r^2,1/(r^2\sin^2\theta))\), raise the index of \(A_i=(r^2,r^2\theta,r^2\sin\theta)\).
    \end{enumerate}

    \textbf{Level D:\,Differential operators}
    \begin{enumerate}
        \item Derive the cylindrical divergence formula from \(\nabla\cdot\mathbf{A}=\frac1{\sqrt g}\partial_i(\sqrt g A^i)\).
        \item Derive the spherical divergence formula from \(\nabla\cdot\mathbf{A}=\frac1{\sqrt g}\partial_i(\sqrt g A^i)\).
        \item Derive the cylindrical scalar Laplacian from the general curvilinear formula.
        \item Derive the spherical scalar Laplacian from the general curvilinear formula.
        \item Compute \(\nabla^2f\) in cylindrical coordinates for \(f=\rho^2+z^2\).
        \item Compute \(\nabla^2f\) in spherical coordinates for \(f=r^2\).
    \end{enumerate}
\end{practicebox}

\ifdefined\OTCOMBINED%
\else
\end{document}
\fi
